%% Copyright 2007-2024 Elsevier Ltd
%% 
%% This file is part of the 'Elsarticle Bundle'.
%% ---------------------------------------------
%% 
%% It may be distributed under the conditions of the LaTeX Project Public
%% License, either version 1.3 of this license or (at your option) any
%% later version.  The latest version of this license is in
%%    http://www.latex-project.org/lppl.txt
%% and version 1.3 or later is part of all distributions of LaTeX
%% version 1999/12/01 or later.
%% 
%% The list of all files belonging to the 'Elsarticle Bundle' is
%% given in the file `manifest.txt'.
%% 
%% Template article for Elsevier's document class `elsarticle'
%% with numbered style bibliographic references
%% SP 2008/03/01
%% $Id: elsarticle-template-num.tex 249 2024-04-06 10:51:24Z rishi $
%%
%\documentclass[preprint,12pt]{elsarticle}

%% Use the option review to obtain double line spacing
%% \documentclass[authoryear,preprint,review,12pt]{elsarticle}

%% Use the options 1p,twocolumn; 3p; 3p,twocolumn; 5p; or 5p,twocolumn
%% for a journal layout:
%% \documentclass[final,1p,times]{elsarticle}
% \documentclass[final,1p,times,twocolumn]{elsarticle}
%% \documentclass[final,3p,times]{elsarticle}

% For Mac
%\documentclass[final,3p,times,twocolumn]{elsarticle} 

% For Windows
\documentclass[final,3p,times,twocolumn, dvipdfmx]{elsarticle}

%% \documentclass[final,5p,times]{elsarticle}

%% For Harv style
%\documentclass[final,5p,times,twocolumn, authoryear]{elsarticle}

%% For including figures, graphicx.sty has been loaded in
%% elsarticle.cls. If you prefer to use the old commands
%% please give \usepackage{epsfig}

%% The amssymb package provides various useful mathematical symbols
\usepackage{amssymb}
%% The amsmath package provides various useful equation environments.
\usepackage{amsmath}
%% The amsthm package provides extended theorem environments
%% \usepackage{amsthm}


\usepackage{subcaption}
\usepackage{here}
\usepackage{bm}

\usepackage{siunitx}
\sisetup{
  detect-all,
  separate-uncertainty = true,
  per-mode = symbol
}



%% The lineno packages adds line numbers. Start line numbering with
%% \begin{linenumbers}, end it with \end{linenumbers}. Or switch it on
%% for the whole article with \linenumbers.
%% \usepackage{lineno}

\journal{AI Thermal Fluids}

\begin{document}

\begin{frontmatter}

%% Title, authors and addresses

%% use the tnoteref command within \title for footnotes;
%% use the tnotetext command for theassociated footnote;
%% use the fnref command within \author or \affiliation for footnotes;
%% use the fntext command for theassociated footnote;
%% use the corref command within \author for corresponding author footnotes;
%% use the cortext command for theassociated footnote;
%% use the ead command for the email address,
%% and the form \ead[url] for the home page:
%% \title{Title\tnoteref{label1}}
%% \tnotetext[label1]{}
%% 
\author[1]{Yoshitaka Ueki\corref{cor1}}
\ead{ueki@rs.tus.ac.jp}
%% \ead[url]{home page}
%\fntext[label2]{}
\cortext[cor1]{Corresponding author}
%% \affiliation{organization={},
%%             addressline={},
%%             city={},
%%             postcode={},
%%             state={},
%%             country={}}
%\fntext[label3]{}

\author[1]{*****} %% Author name
\author[2]{*****} %% Author name

\title{Explainable Convolutional Neural Networks ....}




%% Author affiliation
\affiliation[1]{organization={Department of Applied Electronics, Faculty of Advanced Engineering, Tokyo University of Science},%Department and Organization
            addressline={Niijuku 6-3-1}, 
            city={Katsushika},
            postcode={125-8585}, 
            state={Tokyo},
            country={Japan}}

\affiliation[2]{organization={Oarai Research and Development Institute, Japan Atomic Energy Agency},%Department and Organization
            addressline={Naritacho 4002}, 
            city={Ooarai},
            postcode={311-1393}, 
            state={Ibaraki},
            country={Japan}}

%% Abstract
\begin{abstract}
%% Text of abstract
In this study, a deep-learning-based methodology was developed ...
\end{abstract}

%%Graphical abstract
%\begin{graphicalabstract}
%  \begin{figure}[h]
%    \centering
%    \includegraphics[width=1.5\linewidth]{graphicalabstract.pdf}
%  \end{figure}
%\end{graphicalabstract}

%%Research highlights
\begin{highlights}
\item CNN using acoustic data in TFR detects onset of nucleate boiling accurately
\end{highlights}

%% Keywords
\begin{keyword}
%% keywords here, in the form: keyword \sep keyword

%% PACS codes here, in the form: \PACS code \sep code

%% MSC codes here, in the form: \MSC code \sep code
%% or \MSC[2008] code \sep code (2000 is the default)
Nucleate boiling \sep Acoustic sensing \sep Deep learning \sep Explainable AI \sep Convolutional neural network
\end{keyword}

\end{frontmatter}

%% Add \usepackage{lineno} before \begin{document} and uncomment 
%% following line to enable line numbers
%% \linenumbers

%% main text
%%

%% Use \section commands to start a section
\section{Introduction}
\label{sec1}

Boiling is a highly efficient phenomenon for thermal energy transport ...




\bibliographystyle{elsarticle-num} 
\bibliography{ATE.bib}
%%%%%%%%%%%%%%%%%%%%%%%%%%%%%%%%%%%%%%%%%%%%%%%%%%%

%%%%%%%%%%%%%%%%%%%%%%%%%%%%%%%%%%%%%%%%%%%%%%%%%%%
\end{document}
\endinput
%%
%% End of file `elsarticle-template-num.tex'.
